% EPTCS-style scaffold for a PLACES research-paper submission.
% PLACES 2026 allowed 8 pages for research papers, excluding bibliography and
% appendices. Keep the main body focused and put mechanization details that are
% not needed for the core argument in the appendix.
\documentclass[submission,copyright,creativecommons]{eptcs}
\providecommand{\event}{PLACES 2027}

\usepackage{iftex}
\ifpdf
  \usepackage{underscore}
  \usepackage[T1]{fontenc}
\else
  \usepackage{breakurl}
\fi

\usepackage{amsmath}
\usepackage{amssymb}
\usepackage{amsthm}
\usepackage{mathtools}
\usepackage{placeins}
\usepackage{cleveref}
\usepackage{svg}
\svgsetup{inkscapelatex=false}

\usepackage{tikz}
\usetikzlibrary{decorations.pathmorphing}

\makeatletter
\newcommand{\overdecoration}[2]{%
  \mathpalette{\overdecoration@{#1}}{#2}%
}
\newcommand{\overdecoration@}[3]{%
  \tikz[baseline=(content.base)]{
    \node[
      inner sep=0pt,
      outer sep=0pt
    ] (content) {$#2#3$};

    \draw[#1]
      ([yshift=1pt]content.north west) --
      ([yshift=1pt]content.north east);
  }%
}
\makeatother

\newcommand{\overwiggle}[1]{%
  \overdecoration{
    decorate,
    decoration={
      snake,
      amplitude=.25pt,
      segment length=4pt
    }
  }{#1}%
}

\newcommand{\overdotted}[1]{%
  \overdecoration{
    densely dotted,
    line cap=round,
    line width=.35pt
  }{#1}%
}

\newcommand{\overdashed}[1]{%
  \overdecoration{
    densely dashed,
    line width=.35pt
  }{#1}%
}

\title{Stopping Duality at $\mu$: Yet Another Approach to Recursive Session Types}
\author{Peter Thiemann
\institute{University of Freiburg\\Freiburg, Germany}
\email{thiemann@informatik.uni-freiburg.de}
}

\def\titlerunning{Stopping Duality at $\mu$}
\def\authorrunning{P. Thiemann}

\newcommand{\Agda}{Agda}
\newcommand{\AgdaHen}{\includesvg[height=0.75em]{agda}}
\newcommand{\AgdaLink}[2]{\href{src/#1}{\mbox{\texttt{#2}\,\AgdaHen}}}
\newcommand{\AgdaFile}[2]{\AgdaLink{#1.html}{#2}}
\newcommand{\AgdaName}[3]{\AgdaLink{#1.html\##2}{#3}}
\newcommand{\AgdaMathName}[3]{\href{src/#1.html\##2}{\mbox{#3\,\AgdaHen}}}
\newcommand{\AgdaCounterparts}[1]{%
  \par\smallskip
  \noindent{\footnotesize\emph{Agda counterparts:} #1.}%
  \par\smallskip
}
\newcommand{\Dual}{\mathsf{dual}}
\newcommand{\End}{\mathsf{end}}
\newcommand{\COI}{\mathsf{COI}}
\newcommand{\IND}{\mathsf{IND}}
\newcommand{\sem}[1]{[\![#1]\!]}
\newcommand{\SBH}{S_{\mathsf{BH}}}
\newcommand{\NaiveDual}[1]{\overwiggle{#1}}
\newcommand{\SuspDual}[1]{\overline{#1}}
\newcommand{\EqT}[1]{\mathrel{\approx_T^{#1}}}
\newcommand{\InfDual}[1]{\mathsf{dual}_\infty(#1)}
\newcommand{\ConvT}{\mathrel{\equiv_T}}
\newcommand{\ConvS}{\mathrel{\equiv_S}}
\newcommand{\MatchH}{\mathrel{\sim_H}}
\newcommand{\Obs}{\mathrel{\Downarrow}}
\newcommand{\CoDual}{\mathrel{\perp}}
\newcommand{\StopDual}{\mathrel{\perp_{\mathsf{sd}}}}
\newcommand{\HDual}{\mathsf{dual}_H}
\newcommand{\coindrule}[2]{%
  \vcenter{\offinterlineskip
    \halign{\hfil$\displaystyle##$\hfil\cr
      #1\cr
      \noalign{\kern0.35ex\hrule height 0.35pt
               \kern0.18ex\hrule height 0.35pt\kern0.55ex}
      #2\cr}}}
\newcommand{\coindaxiom}[1]{%
  \vcenter{\offinterlineskip
    \halign{\hfil$\displaystyle##$\hfil\cr
      \noalign{\kern0.35ex\hrule height 0.35pt
               \kern0.18ex\hrule height 0.35pt\kern0.55ex}
      #1\cr}}}
\newtheorem{theorem}{Theorem}
\newtheorem{lemma}[theorem]{Lemma}

\begin{document}
\maketitle

\begin{abstract}
  Most presentations of recursive session types use explicit $\mu$-types with 
  an equi-recursive interpretation.
  Duality for these types is usually defined as a total syntactic operation on
  session types, which has led to problems in the past.
  This paper explores a different design: duality is defined on the observable
  communication structure and stops at $\mu$.  The resulting presentation avoids
  problematic interactions between duality, substitution, and unfolding.
  We give a mechanized development of this stopped-duality approach in \Agda{}
  and relate it to a coinductive ground-truth model.
\end{abstract}

\section{Introduction}
\label{sec:introduction}

Recursive session types are commonly written with explicit $\mu$-binders while
being read equi-recursively: a type $\mu X.S$ is equal to its unfolding
$S[\mu X.S/X]$, which is necessary when advancing to its next communication
action. In binary session types, duality embodies the switch in
perspective from client to server, swapping the direction of all
communications.
The first attempt defines duality as a total syntactic operation on the
syntax, swapping inputs with outputs, selection with branching, and recursively
descending in the body of a $\mu$-type \cite{HondaVK98}.  Write
$\NaiveDual{(-)}$ for this naive syntactic duality; in particular,
$
  \NaiveDual{\mu X.S} = \mu X.\NaiveDual{S}.
  $
  
Naive duality is known to be broken because of the interaction between
duality and unrolling substitution
\cite{LindleyM16,gay.vasconcelos:session-types}. 
To recap the issue, consider a session type given by Bernardi and Hennessy
\cite{BernardiH14,BernardiH16}, which receives a channel of its own type and then repeats:
\begin{align*}
  \SBH &= \mu X.\,?\,\mathsf{chan}(X).X&
  &\text{and its naive syntactic dual}&
  \NaiveDual{\SBH} &= \mu X.\,!\,\mathsf{chan}(X).X.                                      
\end{align*}
Unfolding $\NaiveDual{\SBH}$ and dualizing after unfolding $\SBH$ give
different payloads:
\[
\begin{array}{rcll}
\NaiveDual{\SBH}
  &=& !\,\mathsf{chan}(\NaiveDual{\SBH}).\NaiveDual{\SBH}
  & \text{naive duality, then unfold},\\
\NaiveDual{(?\,\mathsf{chan}(\SBH).\SBH)}
  &=& !\,\mathsf{chan}(\SBH).\NaiveDual{\SBH}
  & \text{unfold first}.
\end{array}
\]
While the continuation types are the same, the first payload is
$\NaiveDual{\SBH}$ but the second is $\SBH$. This example shows that
naive duality is incompatible with unfolding of the recursion: the payload type
must remain the same under duality.

Bernardi and Hennessy \cite{BernardiH14,BernardiH16} as well as the textbook
\cite{gay.vasconcelos:session-types} suggest a solution based on closing the
payload type before applying naive duality. Gay and
others~\cite{DBLP:journals/corr/abs-2004-01322} discuss the design space for
duality operators comparing the approaches known at that time.

This paper proposes another, simple alternative, a \emph{stopped-duality
  operator} $\SuspDual{(-)}$.
This operator acts syntactically on the finite communication
prefix that is visible at 
the top level and deliberately stops at recursive binders.  For the example
above, $\SuspDual{\SBH}$ is not reduced by a syntactic rule for $\mu$.
Instead,
duality can only proceed \emph{after} unfolding $\SBH$ when an observation is required; at that
point the exposed input is dualized to an output.  Recursion is therefore
driven by observation, rather than by a syntactic dual
operator that traverses under $\mu$.

\emph{Contributions}
\begin{itemize}
\item A stopped-duality operator for equi-recursive session types formalized
  using type conversion;
\item an \Agda{} formalization using rewrite rules for the observable duality
  equations; and
\item a mechanized comparison with coinductive duality representing ground truth.
\end{itemize}

\section{Session Types With Stopped Duality}
\label{sec:session-types}

We start with some standard definitions.
Functional types $T,U$ describe payloads; they include function types and
session-typed channels.
Session types $S,R$ contain the usual communication actions, termination,
variables, and recursive binders:
\[
\begin{array}{rcl}
  T,U &::=& \mathbf{1}
       \mid \mathsf{int}
       \mid T \times U
       \mid T \to U
       \mid \mathsf{chan}(S),
       \\[0.5ex]
  S,R &::=& \End
       \mid !T.S
       \mid ?T.S
       \mid \oplus\{\ell_i:S_i\}_{i \in I}
       \mid \&\{\ell_i:S_i\}_{i \in I}
       \mid X
       \mid \mu X.S.
\end{array}
\]
\AgdaCounterparts{indexed syntax
\AgdaName{Types.IND1}{Type}{Type},
\AgdaName{Types.IND1}{SType}{SType}, and
\AgdaName{Types.IND1}{GType}{GType}; coinductive syntax
\AgdaName{Types.COI}{Type}{Type},
\AgdaName{Types.COI}{STypeF}{STypeF}, and
\AgdaName{Types.COI}{SType}{SType}}
We assume that recursive types are contractive, that is, there are no types of
the form $\mu X_1 \dots \mu X_n.X_j$, for some $1\le j\le n$.

Write $\sem{-}$ for the standard equi-recursive interpretation of closed types
as infinite (rational) trees.  Payload constructors and communication heads are
interpreted homomorphically; a recursive type is interpreted by unfolding it, so
$\sem{\mu X.S} = \sem{S[\mu X.S/X]}$.  Because the notation is used only for
closed types, the output is an infinite tree marked with the type constructors,
but without variables $X$ and recursion operators $\mu X.S$.  We use
calligraphic metavariables for such trees: $\mathcal T,\mathcal U$ range over
payload-type trees and $\mathcal S,\mathcal R$ range over session-type trees.
Plain $T,U$ and $S,R$ continue to range over syntactic $\mu$-types.

\subsection{Ground Truth: Relational Duality}
\label{sec:ground-truth}

To compare the upcoming definition of stopped duality in
\cref{sec:stopped-duality} with an extensional account, we use a coinductively
defined ground truth.  As payload types may contain channels, equivalence is staged.
Figure~\ref{fig:ground-duality} first defines equivalence on payload types
inductively, parameterized by a session relation; it then instantiates that
parameter with coinductive session equivalence and defines relational duality.

\begin{figure}[tp]
\footnotesize
\setlength{\arraycolsep}{0.7ex}
\noindent\textbf{1. Inductive type equivalence, parameterized by a session
  relation $\rho$: $\mathcal T \EqT\rho \mathcal U$}
\[
\begin{array}{c}
\dfrac{}{\mathbf{1} \EqT{\rho} \mathbf{1}}
\quad
\dfrac{}{\mathsf{int} \EqT{\rho} \mathsf{int}}
\quad
\dfrac{\mathcal T_1 \EqT{\rho} \mathcal U_1
       \quad \mathcal T_2 \EqT{\rho} \mathcal U_2}
      {\mathcal T_1 \times \mathcal T_2
       \EqT{\rho}
       \mathcal U_1 \times \mathcal U_2}
\quad
\dfrac{\mathcal T_1 \EqT{\rho} \mathcal U_1
       \quad \mathcal T_2 \EqT{\rho} \mathcal U_2}
      {\mathcal T_1 \to \mathcal T_2
       \EqT{\rho}
       \mathcal U_1 \to \mathcal U_2}
\quad
\dfrac{\mathcal S \mathrel{\rho} \mathcal R}
      {\mathsf{chan}(\mathcal S)
       \EqT{\rho}
       \mathsf{chan}(\mathcal R)}
\end{array}
\]
\vspace{-0.5ex}\hrule\vspace{1ex}
\noindent\textbf{2. Coinductive session equivalence: $\mathcal S \approx \mathcal R$}
\[
\begin{array}{c}
\coindrule{\mathcal T \EqT{\approx} \mathcal U
           \quad \mathcal S \approx \mathcal R}
           {!\mathcal T.\mathcal S \approx !\mathcal U.\mathcal R}
\qquad
\coindrule{\mathcal T \EqT{\approx} \mathcal U
           \quad \mathcal S \approx \mathcal R}
           {?\mathcal T.\mathcal S \approx ?\mathcal U.\mathcal R}
\qquad
\coindaxiom{\End \approx \End}
\\[3ex]
\coindrule{\forall i \in I.\ \mathcal S_i \approx \mathcal R_i}
           {\oplus\{\ell_i:\mathcal S_i\}_{i \in I} \approx
            \oplus\{\ell_i:\mathcal R_i\}_{i \in I}}
\qquad
\coindrule{\forall i \in I.\ \mathcal S_i \approx \mathcal R_i}
           {\&\{\ell_i:\mathcal S_i\}_{i \in I} \approx
            \&\{\ell_i:\mathcal R_i\}_{i \in I}}
\end{array}
\]
\vspace{-0.5ex}\hrule\vspace{1ex}
\noindent\textbf{3. Coinductive relational duality: $\mathcal S \CoDual \mathcal R$}
\[
\begin{array}{c}
\coindrule{\mathcal T \EqT{\approx} \mathcal U
           \quad \mathcal S \CoDual \mathcal R}
           {!\mathcal T.\mathcal S \CoDual ?\mathcal U.\mathcal R}
\qquad
\coindrule{\mathcal T \EqT{\approx} \mathcal U
           \quad \mathcal S \CoDual \mathcal R}
           {?\mathcal T.\mathcal S \CoDual !\mathcal U.\mathcal R}
\qquad
\coindaxiom{\End \CoDual \End}
\\[3ex]
\coindrule{\forall i \in I.\ \mathcal S_i \CoDual \mathcal R_i}
           {\oplus\{\ell_i:\mathcal S_i\}_{i \in I} \CoDual
            \&\{\ell_i:\mathcal R_i\}_{i \in I}}
\qquad
\coindrule{\forall i \in I.\ \mathcal S_i \CoDual \mathcal R_i}
           {\&\{\ell_i:\mathcal S_i\}_{i \in I} \CoDual
            \oplus\{\ell_i:\mathcal R_i\}_{i \in I}}
\end{array}
\]
\AgdaCounterparts{coinductive model
\AgdaName{Types.COI}{EquivT}{EquivT},
\AgdaName{Types.COI}{EquivF}{EquivF}, and
\AgdaName{Types.COI}{Equiv}{Equiv}, with the recursive-syntax analogue
\AgdaName{Types.IND1}{EquivT}{EquivT},
\AgdaName{Types.IND1}{EquivG}{EquivG}, and
\AgdaName{Types.IND1}{Equiv}{Equiv}; ground relation
\AgdaName{Types.COI}{DualD}{DualD},
\AgdaName{Types.COI}{DualF}{DualF}, and
\AgdaName{Types.COI}{Dual}{Dual}}
\caption{Equivalence and relational duality in the coinductive model.
Payload equivalence is inductive and parameterized by a session relation.
Session equivalence and relational duality are coinductive; the double-line
rules are read coinductively.}
\label{fig:equivalence}
\label{fig:ground-duality}
\end{figure}

%\FloatBarrier
The coinductive model also admits a dual operator $\InfDual{\mathcal S}$,
which works on session-type trees.  It is defined corecursively by swapping
the visible direction on prefixes, for example:
\[
\begin{array}{rclcrcl}
\InfDual{!\mathcal T.\mathcal S}
  &=& ?\mathcal T.\InfDual{\mathcal S},
&\qquad&
\InfDual{?\mathcal T.\mathcal S}
  &=& !\mathcal T.\InfDual{\mathcal S}.
\end{array}
\]
It also exchanges selection and branching pointwise and maps $\End$ to
$\End$.  Payload types are left unchanged.
Thus $\CoDual$ and $\InfDual{-}$ are relational and functional presentations
of the same coinductive notion, up to session equivalence $\approx$.

\subsection{Stopped Duality}
\label{sec:stopped-duality}

We add stopped duality to the syntax of session types.
\[
\begin{array}{rcl}
S,R &::=& \cdots \mid \SuspDual{S}
\end{array}
\]
\AgdaCounterparts{stopped duality
\AgdaName{DualStopAtMu}{dualS}{dualS} on session terms and
\AgdaName{DualStopAtMu}{dualG}{dualG} on communication heads}
Figure~\ref{fig:conversion} defines the meaning of the operator by type
conversion.
We write $T \ConvT U$ for payload conversion
and $S \ConvS R$ for session conversion.  Both judgments provide inductive
definitions for congruence relations.  Payload conversion calls session
conversion at channel types; session conversion calls payload conversion at
input and output prefixes.
\AgdaCounterparts{syntactic conversion judgments
\AgdaName{Conversion}{ConvT}{ConvT} for payloads and
\AgdaName{Conversion}{ConvS}{ConvS} for sessions}

\begin{figure}[tp]
\footnotesize
\setlength{\arraycolsep}{0.7ex}
\setlength{\abovedisplayskip}{0.5ex}
\setlength{\belowdisplayskip}{0.5ex}
\setlength{\abovedisplayshortskip}{0.5ex}
\setlength{\belowdisplayshortskip}{0.5ex}
\noindent\textbf{1. Type conversion: $T \ConvT U$}
\[
\begin{array}{c}
\dfrac{}{\mathbf{1} \ConvT \mathbf{1}}
\quad
\dfrac{}{\mathsf{int} \ConvT \mathsf{int}}
\quad
\dfrac{T_1 \ConvT U_1 \quad T_2 \ConvT U_2}
      {T_1 \times T_2 \ConvT U_1 \times U_2}
\quad
\dfrac{T_1 \ConvT U_1 \quad T_2 \ConvT U_2}
      {T_1 \to T_2 \ConvT U_1 \to U_2}
\quad
\dfrac{S \ConvS R}
      {\mathsf{chan}(S) \ConvT \mathsf{chan}(R)}
\end{array}
\]
\vspace{-0.5ex}\hrule\vspace{1ex}
\noindent\textbf{2. Session conversion: $S \ConvS R$}
\[
\begin{array}{c}
\dfrac{}{S \ConvS S}
\quad
\dfrac{S \ConvS R}{R \ConvS S}
\quad
\dfrac{S \ConvS R \quad R \ConvS Q}{S \ConvS Q}
\quad
\dfrac{S \ConvS R}{\SuspDual{S} \ConvS \SuspDual{R}}
\quad
\dfrac{}{\mu X.S \ConvS S[\mu X.S/X]}
\\[3ex]
\dfrac{T \ConvT U \quad S \ConvS R}{!T.S \ConvS !U.R}
\quad
\dfrac{T \ConvT U \quad S \ConvS R}{?T.S \ConvS ?U.R}
\qquad
\dfrac{\forall i \in I.\ S_i \ConvS R_i}
      {\oplus\{\ell_i:S_i\}_{i \in I} \ConvS
       \oplus\{\ell_i:R_i\}_{i \in I}}
\quad
\dfrac{\forall i \in I.\ S_i \ConvS R_i}
      {\&\{\ell_i:S_i\}_{i \in I} \ConvS
       \&\{\ell_i:R_i\}_{i \in I}}
\end{array}
\]
\vspace{-0.5ex}\hrule\vspace{1ex}
\noindent\textbf{3. Stopped-duality conversion axioms}
\begin{align*}
  \SuspDual{\End} &\ConvS \End
  & \SuspDual{!T.S} &\ConvS ?T.\SuspDual{S}
  & \SuspDual{?T.S} &\ConvS !T.\SuspDual{S}
  & \SuspDual{\oplus\{\ell_i:S_i\}_{i \in I}}
  &\ConvS \&\{\ell_i:\SuspDual{S_i}\}_{i \in I}
  & \SuspDual{\&\{\ell_i:S_i\}_{i \in I}}
  &\ConvS \oplus\{\ell_i:\SuspDual{S_i}\}_{i \in I}
\end{align*}
\AgdaCounterparts{duality rewrite rules
\AgdaName{DualStopAtMu}{dual-gdd}{dual-gdd},
\AgdaName{DualStopAtMu}{dual-transmit}{dual-transmit},
\AgdaName{DualStopAtMu}{dual-choice}{dual-choice}, and
\AgdaName{DualStopAtMu}{dual-end}{dual-end}; conversion witnesses
\AgdaName{Conversion}{ConvS.conv-dual-transmit}{conv-dual-transmit},
\AgdaName{Conversion}{ConvS.conv-dual-choice}{conv-dual-choice}, and
\AgdaName{Conversion}{ConvS.conv-dual-end}{conv-dual-end}}
\caption{Conversion for stopped-duality session types.  Compartment 3 is the
only computational content of suspended duality.}
\label{fig:conversion}
\end{figure}

%\FloatBarrier
No conversion axiom applies directly to $\SuspDual{\mu X.S}$ or $\SuspDual{X}$.
To obtain the dual of a recursive type, conversion first unfolds by
congruence and then applies a head-duality axiom.  For the
Bernardi-Hennessy example,
\[
\SuspDual{\SBH}
  \ConvS
\SuspDual{?\,\mathsf{chan}(\SBH).\SBH}
  \ConvS
!\,\mathsf{chan}(\SBH).\SuspDual{\SBH},
\]
thus the payload remains $\SBH$. Conversion exposes the
communication structure on demand.

Figure~\ref{fig:duality} defines the observational relation induced by
stopped duality.  It observes the stopped dual $\SuspDual{S}$ and matches the
resulting communication head against the other endpoint.

\begin{figure}[tp]
\footnotesize
\setlength{\arraycolsep}{0.7ex}
\noindent\textbf{1. Sessions in head normal form $H, K$ and head duality: $\HDual(H)=S$}
\begin{align*}
H,K &::= \End \mid !T.S \mid ?T.S
      \mid \oplus\{\ell_i:S_i\}_{i \in I}
      \mid \&\{\ell_i:S_i\}_{i \in I}
\end{align*}
\begin{align*}
\HDual(\End) &= \End &
\HDual(!T.S) &= ?T.\SuspDual{S} &
\HDual(\oplus\{\ell_i:S_i\}_{i \in I}) &=
  \&\{\ell_i:\SuspDual{S_i}\}_{i \in I} \\
&&
\HDual(?T.S) &= !T.\SuspDual{S} &
\HDual(\&\{\ell_i:S_i\}_{i \in I}) &=
  \oplus\{\ell_i:\SuspDual{S_i}\}_{i \in I}
\end{align*}
\vspace{-0.5ex}\hrule\vspace{1ex}
\noindent\textbf{2. Head matching of stopped-duality terms: $S \MatchH H$}
\[
\begin{array}{c}
\dfrac{}{H \Obs H}
\qquad
\dfrac{S_0[\mu X.S_0/X] \Obs H}
      {\mu X.S_0 \Obs H}
\qquad
\dfrac{S \Obs H}
      {\SuspDual{S} \Obs \HDual(H)}
\end{array}
\]
\vspace{-0.5ex}\hrule\vspace{1ex}
\noindent\textbf{3. Observational duality induced by stopped duality: $S
  \StopDual R$}
\[
\begin{array}{c}
\dfrac{}{\End \MatchH \End}
\qquad
\dfrac{T \EqT{\approx} U \quad S \StopDual R}
      {!T.S \MatchH !U.R}
\qquad
\dfrac{T \EqT{\approx} U \quad S \StopDual R}
      {?T.S \MatchH ?U.R}
\\[3ex]
\dfrac{\forall i \in I.\ S_i \StopDual R_i}
      {\oplus\{\ell_i:S_i\}_{i \in I} \MatchH
       \oplus\{\ell_i:R_i\}_{i \in I}}
\qquad
\dfrac{\forall i \in I.\ S_i \StopDual R_i}
      {\&\{\ell_i:S_i\}_{i \in I} \MatchH
       \&\{\ell_i:R_i\}_{i \in I}}
\\[3ex]
\coindrule{\SuspDual{S} \Obs H
           \quad R \Obs K
          \quad H \MatchH K}
          {S \StopDual R}
\end{array}
\]
\AgdaCounterparts{unfolding
\AgdaName{Types.IND1}{unfold}{unfold} and stopped duality
\AgdaName{DualStopAtMu}{StopDualG}{StopDualG} and
\AgdaName{DualStopAtMu}{StopDual}{StopDual}}
\caption{Stopped duality.  Observation unfolds recursive types before
$\HDual$ switches the visible action.  The head-matching relation $\MatchH$
compares payloads with type equivalence from Figure~\ref{fig:equivalence} and
continues coinductively through $\StopDual$.}
\label{fig:duality}
\end{figure}

%\FloatBarrier
\subsection{Comparison}
\label{sec:comparison}

Conversion and equivalence have different roles.  Figure~\ref{fig:conversion}
defines conversion $\ConvT,\ConvS$ as a syntactic equality judgment for the
extended grammar with suspended duals; it includes $\mu$-unfolding and the
stopped-duality axioms and is intended for type checking.  The parameterized
equivalence $\EqT{\rho}$ and its coinductive session instance $\approx$ are
extensional relations used to compare behavior after interpretation.  Their
connection is semantic soundness.
\begin{theorem}[conversion soundness]
\label{thm:conversion-soundness}
For all closed extended payload types $T,U$ and closed extended session types
$S,R$,
\begin{align*}
T \ConvT U
  &\Longrightarrow \sem{T} \EqT{\approx} \sem{U},
  &
S \ConvS R
  &\Longrightarrow \sem{S} \approx \sem{R}.
\end{align*}
\end{theorem}
\AgdaCounterparts{conversion soundness
\AgdaName{Conversion}{convT-sound}{convT-sound} and
\AgdaName{Conversion}{convS-sound}{convS-sound}; duality congruence
\AgdaName{Conversion}{dual-cong}{dual-cong}; unrolling adequacy
\AgdaName{Conversion}{unroll-sound}{unroll-sound}; definitional adequacy
\AgdaName{Conversion}{dual-transmit-adequacy}{dual-transmit-adequacy},
\AgdaName{Conversion}{dual-choice-adequacy}{dual-choice-adequacy}, and
\AgdaName{Conversion}{dual-end-adequacy}{dual-end-adequacy}}
The inductive conversion judgment is not complete for regular-tree
equality.  For example, coinductive stopped duality can relate
$S_1 = \mu X.?\,\mathsf{Unit}.X$ to
$S_2 = \mu X.! \,\mathsf{Unit}.! \,\mathsf{Unit}.X$, but $\SuspDual{S_1}$ is not
convertible to $S_2$: the observation repeats
coinductively after one step on the left and two on the right.
\AgdaCounterparts{counterexample
\AgdaName{Examples}{ConversionCounterexample.recv-loop}{recv-loop} and
\AgdaName{Examples}{ConversionCounterexample.send-loop2}{send-loop2};
\AgdaName{Examples}{ConversionCounterexample.dual-recv-loop-step}{dual-step}
and \AgdaName{Examples}{ConversionCounterexample.send-loop2-step}{send-step};
\AgdaName{Examples}{ConversionCounterexample.recv-loop-dual-send-loop2}
{sd-proof} and
\AgdaName{Examples}{ConversionCounterexample.recv-loop-dual-send-loop2-coi}
{coi-proof}}


The stopped relation, the coinductive ground truth, and the stopped-duality
operator coincide on closed session types.
\begin{theorem}[master comparison]
\label{thm:master-comparison}
For all closed session types $S$ and $R$, the following statements are
equivalent:
\[
\begin{array}{ll}
\mathrm{(i)}   & S \StopDual R,\\[0.35ex]
\mathrm{(ii)}  & \sem{S} \CoDual \sem{R},\\[0.35ex]
\mathrm{(iii)} & \InfDual{\sem{S}} \approx \sem{R},\\[0.35ex]
\mathrm{(iv)}  & \sem{\SuspDual{S}} \approx \sem{R}.
\end{array}
\]
\end{theorem}
\AgdaCounterparts{the fields
\AgdaName{ConversionStopDualSoundness}{MasterComparison.sd-coi}{sd-coi},
\AgdaName{ConversionStopDualSoundness}{MasterComparison.sd-functional}{sd-functional},
and
\AgdaName{ConversionStopDualSoundness}{MasterComparison.sd-stopped-coi}{sd-stopped-coi}
of the record
\AgdaName{ConversionStopDualSoundness}{MasterComparison}{MasterComparison}, with
proof \AgdaName{ConversionStopDualSoundness}{master-comparison}{master-comparison}}
The equivalence between (i) and (ii) calibrates observational stopped duality
against the coinductive relation.  Statement (iii) replaces relational
duality by the functional coinductive dual.  Statement (iv) says that the
stopped operator denotes the same coinductive dual.

%\FloatBarrier
\section{Examples}
\label{sec:examples}

The module \AgdaFile{Examples}{Examples} collects small communication examples
and the legacy counterexample that motivated the design.  These examples are
intended to serve two roles in the paper: compact sanity checks for the
proposed operation, such as \AgdaName{Examples}{stopped-send-end}{stopped-send-end}
and \AgdaName{Examples}{stopped-bh}{stopped-bh}, and evidence that the stopped
presentation avoids an otherwise awkward duality-substitution interaction,
captured by \AgdaName{Examples}{SubstitutionDualCounterexample.bh-not-involutive}{bh-not-involutive}.



\section{Mechanization}
\label{sec:mechanization}

The presentation of types and session types in the mechanization relies on de
Bruijn indices to represent variables, it separates session and branching
structure, and enforces contractiveness of recursive types by construction.
Identifiers marked by \AgdaHen{} link to generated \Agda{} HTML.  The links
resolve when the accompanying archive is unpacked next to the PDF, creating a
sibling directory \texttt{src/}.

The mechanization represents the stopped-duality conversion axioms as rewrite
rules \cite{cockx:LIPIcs.TYPES.2019.2,10.1145/3434341} for \AgdaName{DualStopAtMu}{dualS}{dualS} and
\AgdaName{DualStopAtMu}{dualG}{dualG}.  The formal soundness proof proceeds by
induction on \AgdaName{Conversion}{ConvT}{ConvT} and
\AgdaName{Conversion}{ConvS}{ConvS}.  The stopped-duality cases are discharged
by definitional equality before the semantic comparison is applied; the
$\mu$-rule uses the adequacy of syntactic unrolling for the stack
interpretation.  The remaining difference is structural: conversion is an
inductive syntactic judgment with explicit rules for suspended duals, whereas
equivalence is coinductive, parameterized by a session relation, and stated
after interpretation.


Figure~\ref{fig:ground-duality} has a direct counterpart in
\AgdaFile{Types.COI}{Types.COI}: \AgdaName{Types.COI}{DualD}{DualD} relates
opposite directions, \AgdaName{Types.COI}{DualF}{DualF} relates one layer of
communication structure, and the coinductive record
\AgdaName{Types.COI}{Dual}{Dual} defines the relation $\CoDual$.  The same
module implements the auxiliary relations from Figure~\ref{fig:equivalence}
as \AgdaName{Types.COI}{EquivT}{EquivT},
\AgdaName{Types.COI}{EquivF}{EquivF}, and
\AgdaName{Types.COI}{Equiv}{Equiv}.

The module \AgdaFile{DualStopAtMu}{DualStopAtMu} postulates the stopped
operators \AgdaName{DualStopAtMu}{dualS}{dualS} and
\AgdaName{DualStopAtMu}{dualG}{dualG} and gives their observable equations as
rewrite rules for the \AgdaName{DualStopAtMu}{dual-transmit}{transmit},
\AgdaName{DualStopAtMu}{dual-choice}{choice}, and
\AgdaName{DualStopAtMu}{dual-end}{end} cases.  Additional force equations for
\AgdaName{DualStopAtMu}{dual-rec-force}{rec} and
\AgdaName{DualStopAtMu}{dual-var-force}{var} expose the behavior at recursive
binders and variables after interpretation in the coinductive model.
The theorem \AgdaName{DualStopAtMu}{ground}{ground} supplies the equivalence
between statements (iii) and (iv) in Theorem~\ref{thm:master-comparison}.

The module \AgdaFile{Conversion}{Conversion} gives the syntactic conversion
judgments \AgdaName{Conversion}{ConvT}{ConvT} and
\AgdaName{Conversion}{ConvS}{ConvS}.  Its soundness proof is split into
\AgdaName{Conversion}{convT-sound}{convT-sound} for payload conversion and
\AgdaName{Conversion}{convS-sound}{convS-sound} for session conversion.  The
stopped-duality conversion axioms are also exposed as definitional adequacy
lemmas such as
\AgdaName{Conversion}{dual-transmit-adequacy}{dual-transmit-adequacy}.

The module \AgdaFile{DualStopAtMu}{DualStopAtMu} also packages the
observational relation from
Figure~\ref{fig:duality} as the separately named coinductive record
\AgdaName{DualStopAtMu}{StopDual}{StopDual},
with operator notation corresponding to $\StopDual$.  Its observation field
unfolds both closed endpoints and returns one
\AgdaName{DualStopAtMu}{StopDualG}{StopDualG} layer.  The constructors
\AgdaName{DualStopAtMu}{StopDualG.sd-transmit}{sd-transmit},
\AgdaName{DualStopAtMu}{StopDualG.sd-choice}{sd-choice}, and
\AgdaName{DualStopAtMu}{StopDualG.sd-end}{sd-end} fuse the observation and
head-matching parts of the figure: they require opposite visible directions,
equivalent payload types, and recursively
\AgdaName{DualStopAtMu}{StopDual}{StopDual}-related continuations.  It also
provides symmetry
\AgdaMathName{DualStopAtMu}{2574}{$\mathsf{symm}$} and inversion
\AgdaMathName{DualStopAtMu}{4133}{$\mathsf{inv}$} lemmas.

The Agda file \AgdaFile{StopDualSoundness}{SDS} mechanizes the equivalence
between statements (i) and (ii) in Theorem~\ref{thm:master-comparison}.
The main technical bridge,
\AgdaName{StopDualSoundness}{unfold-sound}{unfold-sound}, identifies
syntactic one-step unfolding with forcing the stack interpretation.  The
one-step forms are \AgdaName{StopDualSoundness}{sd-soundG}{sd-soundG} and
\AgdaName{StopDualSoundness}{sd-completeG}{sd-completeG}.  The soundness
direction is organized by the auxiliary coalgebras
\AgdaName{StopDualSoundness}{EqRel}{EqRel} and
\AgdaName{StopDualSoundness}{SdRel}{SdRel}; these relations close the proof
under coinductive equivalence and duality while allowing intermediate
equivalence transports.  The completeness direction uses
\AgdaName{StopDualSoundness}{EqCompleteRel}{EqCompleteRel} and
\AgdaName{StopDualSoundness}{SdCompleteRel}{SdCompleteRel} to reflect the
coinductive observations back to the syntactic relation.  Its payload case
depends on \AgdaName{StopDualSoundness}{equiv-completeT}{equiv-completeT},
because stopped duality stores payload comparison in the inductive syntax
whereas the ground-truth relation produces payload equivalence after
interpretation.

The comparison file \AgdaFile{ConversionStopDualSoundness}{CSDS} packages the
link between conversion and observational stopped duality.  The theorem
\AgdaName{ConversionStopDualSoundness}{conv-dual-sound-sd}{conv-dual-sound-sd}
derives $S \StopDual R$ from a conversion proof
$\SuspDual{S} \ConvS R$.  The formalization also proves a closed-syntax
reflection result, recorded as
\AgdaName{ConversionStopDualSoundness}{MasterComparison.sd-stopped-syntax}{sd-stopped-syntax}.
It combines
\AgdaName{StopDualSoundness}{sd-sound}{sd-sound}, the coinductive
functional/relational duality lemma
\AgdaName{Types.COI}{dual-completeS}{dual-completeS}, the ground theorem
\AgdaName{DualStopAtMu}{ground}{ground}, soundness and reflection of
syntactic equivalence through
\AgdaName{StopDualSoundness}{equiv-soundS}{equiv-soundS} and
\AgdaName{StopDualSoundness}{equiv-completeS}{equiv-completeS}, and transport
of coinductive duality along equivalence.  The
\AgdaName{ConversionStopDualSoundness}{MasterComparison}{record} and
\AgdaName{ConversionStopDualSoundness}{master-comparison}{proof}
package the components of Theorem~\ref{thm:master-comparison} together with
this extra syntactic component.

The mechanization also contains full renamings and simultaneous substitutions
for the de Bruijn representation in \AgdaFile{Types.IND1}{Types.IND1}.
Renamings are represented by \AgdaName{Types.IND1}{Ren}{Ren} and implemented
by \AgdaName{Types.IND1}{renameS}{renameS},
\AgdaName{Types.IND1}{renameG}{renameG}, and
\AgdaName{Types.IND1}{renameT}{renameT}; simultaneous substitutions are
represented by \AgdaName{Types.IND1}{Sub}{Sub} and implemented by
\AgdaName{Types.IND1}{substS}{substS},
\AgdaName{Types.IND1}{substG}{substG}, and
\AgdaName{Types.IND1}{substT}{substT}.  Every type representation now
contains the arrow constructor \texttt{TFun}; the constructor is covered by
renaming, substitution, embeddings between representations, type equivalence
and its proofs, message closure, and congruence contexts.


\section{Related Work}
\label{sec:related-work}

Binary session types originate in the work of Honda and collaborators
\cite{DBLP:conf/concur/Honda93,TakeuchiHK94,HondaVK98}.  In this setting,
duality is the operation that turns the server view of a endpoint into the
client view.  Later work on binary-session subtyping 
\cite{GayH99,gay-hole-subtyping,Vasconcelos12} refines the behavioral
preorder on endpoints, but still relies on compatibility with
duality.  The present paper is orthogonal to the design of a process calculus
or subtyping system: it isolates the definition of duality for recursive
session types with higher-order channel payloads.

The difficulty addressed here is the interaction between syntactic duality and
equi-recursive unfolding.  The traditional definition descends through a
recursive binder, for example by setting
$\NaiveDual{\mu X.S} = \mu X.\NaiveDual{S}$ \cite{HondaVK98}.  Bernardi and
Hennessy exhibited the counterexample used in the introduction in their work on
higher-order contracts for session types \cite{BernardiH14,BernardiH16}.  The
same issue is discussed by Lindley and Morris in the setting of structurally
recursive session programming \cite{LindleyM16}, where recursion variables are
annotated with polarities.  Bernardi and Hennessy and the recent textbook of
Gay and Vasconcelos \cite{gay.vasconcelos:session-types} take a different
route: they avoid dualizing open occurrences in channel payloads, typically by
closing the type before duality is applied.

Gay, Thiemann, and Vasconcelos surveyed this design space in ``Duality of
Session Types: The Final Cut'' \cite{DBLP:journals/corr/abs-2004-01322}.  That
paper compares syntactic duality operators and duality relations for recursive
session types and identifies formulations that avoid the Bernardi-Hennessy
failure.  Stopped duality gives a different presentation of the same problem:
instead of repairing the traversal under $\mu$, it removes that traversal from
the duality operator.  Duality is reduced only at observable communication
heads, while recursive unfolding is handled by conversion and by the
coinductive observation relation.  The comparison theorem in
\cref{sec:comparison} shows that this presentation agrees with the
coinductive ground truth.

There is also related work on defining duality directly as a relation.
Bernardi, Dardha, Gay, and Kouzapas study duality relations for session types
\cite{BernardiDGK14}.  The coinductive relation $\CoDual$ used here is in that
spirit: it is a relational specification of peer behavior over infinite
session trees.  The stopped-duality operator provides a simple approach that we
prove correct with respect to the relation.

The mechanized part follows standard representations of binding with de Bruijn
indices and simultaneous substitutions, but its proof obligations are shaped
by equi-recursive interpretation.  The Agda formalization \cite{agda} uses
rewrite rules \cite{cockx:LIPIcs.TYPES.2019.2,10.1145/3434341} for the computational cases of stopped duality and coinductive
records for the infinite-tree model.  Thus, we obtain executable conversion
principles for the finite syntax and a mechanically checked bridge to the
coinductive semantics.

\section{Conclusion}
\label{sec:conclusion}

Stopped duality is a new alternative presentation of duality for equi-recursive
session types with higher-order channels.  Its computation is confined
to observable communication heads; recursive types are unfolded through
conversion and observation.  Stopped duality is equivalent to other sound
definitions and it gives rise to an elegant mechanization in a proof assistant
like Agda that supports rewriting.

The Agda development validates stopped duality against a coinductive
ground-truth model.  The master comparison theorem relates observational
stopped duality, relational coinductive duality, the functional coinductive
dual, and the interpretation of stopped duality.  Conversion soundness is a
foundation for type-checking.

\bibliographystyle{eptcs}
\bibliography{references}

\appendix

\section{Mechanization Map}

The current development consists of the following files:
\begin{itemize}
\item \AgdaFile{Types.Direction}{Direction}: send/receive
  directions and their involution;
\item \AgdaFile{Types.IND1}{IND1}: indexed recursive syntax,
  renamings, substitutions, weakening, and syntactic equivalence;
\item \AgdaFile{Types.COI}{COI}: coinductive session trees,
  equivalence, relational duality, and functional duality;
\item \AgdaFile{Types.Tail1}{Tail1}: tail-recursive syntax used by
  the stack interpretation and message-closure examples;
\item \AgdaFile{DualTail1}{DualTail1}: stack-based interpretation of
  recursive syntax into the coinductive model;
\item \AgdaFile{DualStopAtMu}{DualStopAtMu}: stopped duality,
  confluence-checked rewrite rules, observational stopped duality, and the
  ground comparison with functional coinductive duality;
\item \AgdaFile{Conversion}{Conversion}: syntactic conversion judgments,
  unrolling adequacy, and conversion soundness;
\item \AgdaFile{StopDualSoundness}{StopDualSoundness}: soundness and
  completeness of observational stopped duality with respect to coinductive
  relational duality;
\item \AgdaFile{ConversionStopDualSoundness}{ConversionStopDualSoundness}:
  conversion, stopped duality, syntactic equivalence, and the master
  comparison theorem;
\item \AgdaFile{ConversionCompletenessCounterexample}
  {conversion counterexample}: the period-collapse example
  showing that inductive conversion does not capture all coinductive
  regular-tree behavior;
\item \AgdaFile{MessageClosure}{MessageClosure}: translation from open
  recursive syntax to tail syntax with closed channel payloads;
\item \AgdaFile{Examples}{Examples}: examples and counterexamples used by
  the paper; and
\item \AgdaFile{Auxiliary.Extensionality}{Extensionality}: the
  functional extensionality postulate used by syntax-manipulation lemmas.
\end{itemize}

The repository also retains older or out-of-focus developments:
\texttt{Types/IND.agda}, \texttt{Duality.agda}, \texttt{DualRel.agda},
\texttt{STypeCongruence.agda}, and the modules under \texttt{OutOfFocus/}.
They are not used by the results stated in this paper.

\end{document}
